% !TeX root = ../常微分方程笔记.tex
% 由主文档通过 \input 载入；请编译主文档。
\setcounter{chapter}{7} % 保留原书章号。
\chapter{自治系统与平面定性理论}
\label{chap:ode-phase-plane}

本章设 $x'=f(x)$，$f$ 在开区域 $D\subseteq\mathbb R^d$ 上为 $C^1$。即使没有显式解，也可以由向量场判断解的方向、可能的极限与周期运动。

\section{相空间、轨道与局部流}
状态空间称相空间（phase space），解 $x(t)$ 在其中的像称轨道（orbit）。记初值为 $x_0$ 的解为 $\varphi_t(x_0)$，称局部流（local flow）。自治性和唯一性给出
\[
 \varphi_0(x_0)=x_0,\qquad
 \varphi_{t+s}(x_0)=\varphi_t(\varphi_s(x_0)),
\]
其中有关解都必须存在。若两条解轨道相交，将相交时刻对齐后可由唯一性知它们属于同一轨道；不同初始时刻只是同一运动的不同参数起点。

若解在某个 $T>0$ 后回到同一状态，则唯一性使其重复原来的运动。非平衡周期解的轨道是闭曲线；平衡点本身则是一点轨道。

\section{标量自治方程与平衡点}
对 $y'=g(y)$，在不含零点的区间内 $g$ 符号固定，故非平衡解严格单调。实轴上标出 $g=0$ 的点并在各区间画出方向，得到相线（phase line）。若全局正向解有界且单调，它存在有限极限 $\ell$；若 $\ell$ 仍在 $g$ 的定义域内，则必有 $g(\ell)=0$，否则连续性使 $y'$ 在充分大时间后保持远离零，矛盾。

在局部唯一性条件下，非平衡解不能有限时间到达平衡点：若到达，则以该时刻为初值向前、向后应用唯一性，将迫使它与常值解重合。因而 $y'=y(1-y)$ 的正解从两侧趋近 $1$，却不在有限时间越过 $1$。

\section{平面线性系统的相图}
对实 $2\times2$ 矩阵 $A$，记 $\tau=\operatorname{tr}A$、$\Delta=\det A$，特征方程为 $\lambda^2-\tau\lambda+\Delta=0$。相图（phase portrait）汇集区域内不同初值的轨道。
\begin{itemize}
 \item $\Delta<0$：两个实根异号，原点为鞍点（saddle），一条方向收缩、另一条方向伸长。
 \item $\Delta>0$、$\tau^2>4\Delta$：两个不同实根同号，原点为结点（node）；$\tau<0$ 吸引，$\tau>0$ 排斥。
 \item $\tau^2<4\Delta$、$\tau\ne0$：复根实部非零，原点为焦点（focus），按 $\tau$ 的符号向内或向外旋转。
 \item $\tau=0$、$\Delta>0$：纯虚特征值，原点为中心（center），非零轨道是闭合椭圆。
\end{itemize}
重实根情形还须检查是否有两个独立特征向量；$\Delta=0$ 有零特征值，不属于上述非退化分类。旋转方向应由实际向量场判断，不能只由特征值决定。

\section{非线性平面系统的相图构造}
设 $u'=P(u,v)$、$v'=Q(u,v)$。零增长曲线（nullcline）$P=0$、$Q=0$ 把平面划分为分量符号固定的区域：$P>0$ 向右、$Q>0$ 向上，其余类推。两类曲线交点为平衡点，但某一条零增长曲线通常不是轨道。

例如 $u'=u(1-u-v)$、$v'=v(2-u-v)$ 在第一象限中以直线 $u+v=1,2$ 分区。坐标轴因右端含相应变量而保持不变，局部唯一性又阻止正解穿越坐标轴。这样的正向不变集（positively invariant set）帮助限定长期行为。平衡点附近可计算 Jacobian 矩阵作线性化，其严格适用条件见第\ref{chap:ode-stability}章。

\section{首次积分、能量与周期运动}
若 $H\in C^1(D)$ 且 $\nabla H\cdot f=0$，则由链式法则 $\mathrm dH(x(t))/\mathrm dt=0$，故 $H$ 是首次积分。轨道位于其等值集内，但一个等值集可以包含多条轨道。

对无阻尼单摆 $\theta''+\omega^2\sin\theta=0$，$\omega>0$，令 $v=\theta'$，有能量
\[
 H(\theta,v)=\frac{v^2}{2}+\omega^2(1-\cos\theta).
\]
在一个势阱内，$0<E<2\omega^2$ 的能量曲线为围绕最低点的闭合曲线，且向量场处处非零，故对应周期摆动。理由是该正则闭曲线紧，运动速度的模有正下界，一次绕行所需时间有限。$E=2\omega^2$ 的曲线经过不稳定平衡点，不能由同一论证推出有限周期。$E>2\omega^2$ 对应持续旋转；若将角度视为实数，其轨道不闭合，若模 $2\pi$ 识别角度，则应在圆柱相空间中理解。

\section{周期轨道与极限环（选学）}
孤立的周期轨道称极限环（limit cycle），中心周围的一族闭轨不满足孤立性。平面系统
\[
 u'=(1-u^2-v^2)u-v,\qquad v'=u+(1-u^2-v^2)v
\]
在极坐标中满足 $r'=r(1-r^2)$、$\theta'=1$。所以 $r=1$ 是吸引的极限环，非零邻近轨道渐近靠近它而不与之相交。

\begin{proposition}[Bendixson 判据]
在单连通开区域中，若 $P,Q\in C^1$ 且 $P_u+Q_v$ 处处严格同号，则区域内不存在非平衡周期轨道。
\end{proposition}
\begin{proof}
若存在周期轨道，它是简单闭曲线，围成的有界区域也在给定单连通区域内。沿轨道向量场与边界相切，通量为零；Green 公式却把此通量写为 $\iint(P_u+Q_v)\,\mathrm du\,\mathrm dv$，其严格同号，矛盾。
\end{proof}

对正向全局轨道，定义 $\omega$ 极限集（omega-limit set）为所有 $x(t_n)$ 的极限点，其中 $t_n\to\infty$。Poincaré--Bendixson 定理的基础版本是：平面 $C^1$ 系统的一条正向轨道若保持在定义域内部的紧集中，且其 $\omega$ 极限集不含平衡点，则该极限集是一条周期轨道。这里引用该定理而省略平面拓扑证明，见\cite[第7章]{teschl-ode}。维数、紧性和无平衡点条件均不可省略。

\section{习题}
\begin{enumerate}
 \item\label{ex:ode-phase} 分析 $y'=y(1-y)(y-2)$ 的相线及各平衡点附近的方向。
 \item\label{ex:ode-spiral} 判断 $A=\begin{pmatrix}-1&-2\\2&-1\end{pmatrix}$ 的相图类型与旋转方向。
 \item 对上述极限环系统，证明任意 $r(0)>0$ 的解满足 $r(t)\to1$。
 \item 证明 $u'=-u+\sin v$、$v'=-v+\sin u$ 在平面上没有非平衡周期轨道。
\end{enumerate}

